\begin{center}
{\Huge \scshape Qiong Liu, Ph.D.} \\ \vspace{4pt}
\small
\faPhone\ 475-280-1364 \quad
\href{mailto:qiongliu.work@gmail.com}{qiongliu.work@gmail.com} \quad
\href{https://www.linkedin.com/in/qiong-l-529b23196}{LinkedIn} \quad
\href{https://scholar.google.com/citations?user=SvnpHCkAAAAJ&hl=en}{Google Scholar}

Authorized to work in the U.S. without sponsorship
\end{center}

\section{Professional Summary}
\textbf{Computer Vision / AI Researcher} with experience developing deep learning methods for image restoration, motion correction, denoising, and spatiotemporal modeling in real-world imaging systems. Strong background in \textbf{representation learning, self-supervised learning, deformable modeling, and robust evaluation} with first-author publications and patent-pending work. Experienced in translating research ideas into deployable algorithms for clinically meaningful image quality improvement.

\section{Core Expertise}
\textbf{Computer Vision \& ML:} Image restoration, denoising, deformable registration, representation learning, self-supervised learning, diffusion models, transformers \\
\textbf{Modeling:} Spatiotemporal modeling, motion modeling, signal decomposition, optimization, uncertainty-aware evaluation \\
\textbf{Programming:} Python, PyTorch, NumPy, SciPy, MATLAB \\
\textbf{Research Strengths:} Problem formulation, algorithm design, evaluation, cross-functional collaboration

\section{Experience}
\textbf{Canon Medical Research USA} \hfill Apr 2025 -- Present \\
\textit{AI Scientist (Reconstruction), Vernon Hills, IL}
\begin{itemize}[leftmargin=*]
\item Developed deep learning algorithms for \textbf{motion correction and image quality improvement} in dynamic imaging under noisy and incomplete data conditions
\item Designed \textbf{spatiotemporal modeling} methods to improve robustness and consistency across sequential image frames
\item Built training and evaluation pipelines for real-world imaging data, integrating preprocessing, model training, inference, and quantitative assessment
\item Collaborated with research and product teams to translate algorithm prototypes into deployable vision solutions
\end{itemize}

\textbf{Canon Medical Research USA} \hfill Jun 2023 -- May 2024 \\
\textit{Research Scientist Intern, Vernon Hills, IL}
\begin{itemize}[leftmargin=*]
\item Developed deep learning methods for \textbf{image denoising and low-signal restoration}, improving image fidelity while preserving critical structures
\item Designed evaluation frameworks for robustness, quantitative consistency, and generalization across heterogeneous datasets
\item Presented technical findings to guide model design and downstream product decisions
\end{itemize}

\textbf{United Imaging Healthcare} \hfill Jul 2019 -- Aug 2019 \\
\textit{Student Intern, Wuhan, China}
\begin{itemize}[leftmargin=*]
\item Supported development and validation of engineering systems for clinical imaging applications
\end{itemize}

\section{Patents}
\textbf{Edge-Guided Deformation Field Fusion for Motion Correction} \hfill 2026 \\
Patent Pending, \textit{Lead Inventor}
\begin{itemize}[leftmargin=*]
\item Proposed a spatially adaptive deformable modeling framework to improve alignment accuracy and robustness in challenging image regions
\end{itemize}

\textbf{Automatic Data-Driven Multi-Signal Decomposition} \hfill 2025 \\
Patent Pending, \textit{Lead Inventor}
\begin{itemize}[leftmargin=*]
\item Developed an AI-based framework to separate overlapping temporal signals from dynamic image sequences
\end{itemize}

\section{Selected Research Projects}

\textbf{Physiology-Guided Motion Correction for Dynamic PET}
\begin{itemize}[leftmargin=*]
\item Developed learning-based motion correction methods for dynamic PET, addressing failure modes caused by \textbf{noise disparity, mixed Gaussian/Poisson noise, and tracer-dependent uptake heterogeneity} in real patient images
\item Identified limitations of conventional network constraints based on global smoothness priors, which were often insufficient for anatomically ambiguous and low-count regions
\item Proposed a \textbf{physiology-guided mask} to provide spatial guidance during network training, improving robustness in regions with challenging tracer distribution and low image contrast
\item Introduced \textbf{spatially varying regularization} in place of a global smoothness constraint, enabling better balance between local flexibility and stable deformation estimation
\end{itemize}

\textbf{Kinetics-Guided Self-Supervised Denoising for Dynamic PET}
\begin{itemize}[leftmargin=*]
\item Developed \textbf{Patlak-Guided Self-Supervised Denoising Network (PG-SSDNet)} for dynamic PET denoising to improve the quality and quantitative reliability of parametric Ki imaging
\item Incorporated a \textbf{Patlak consistency loss} into a multi-frame self-supervised framework to preserve tracer kinetic behavior across dynamic frames
\item Demonstrated that \textbf{integrating physics/kinetic modeling into AI training} improves denoising performance while maintaining physiologically meaningful quantitative endpoints
\item Evaluated the method on whole-body \textsuperscript{18}F-FDG dynamic PET datasets and showed improved Patlak fitting consistency compared with supervised and unguided baselines
\end{itemize}

\textbf{Deep Image Prior for Dynamic PET Denoising and Quantification}
\begin{itemize}[leftmargin=*]
\item Developed population-based deep image prior methods for dynamic PET denoising, with emphasis on preserving downstream \textbf{parametric quantification accuracy}
\item Investigated how \textbf{training data composition and population priors} affect denoising behavior and quantitative outcome, showing that dataset design materially influences model performance
\item Analyzed CNN training behavior and evolving noise patterns to better understand how deep models learn image structure in low-count dynamic PET
\item Demonstrated that, although deep learning models are often treated as black boxes, their behavior can still be \textbf{systematically interpreted and leveraged} to improve both image quality and quantitative reliability
\end{itemize}

\section{Education}
\textbf{Yale University} \hfill Aug 2020 -- May 2025 \\
Ph.D. in Biomedical Engineering \\
Thesis: Improving Image Quality and Quantification Accuracy for Static and Dynamic PET

\textbf{Huazhong University of Science and Technology} \hfill Sep 2016 -- Jun 2020 \\
B.S. in Biomedical Engineering, Minor in Computer Science \\
GPA: 3.94/4.00

\section{Selected Publications}
\textbf{First Author}
\begin{itemize}[leftmargin=*]
\item Liu Q., et al., \textit{Parametric cardiac imaging with $^{18}$F-flutemetamol PET to evaluate the impact of tafamidis in patients with transthyretin cardiac amyloidosis}, \textbf{Journal of Nuclear Medicine}, 2026
\item Liu Q., et al., \textit{Patlak-guided self-supervised learning for dynamic PET denoising}, \textbf{IEEE Transactions on Radiation and Plasma Medical Sciences}, 2026
\item Liu Q., et al., \textit{Population-based deep image prior for dynamic PET denoising: a data-driven approach to improve parametric quantification}, \textbf{Medical Image Analysis}, 2024
\item Liu Q., et al., \textit{A personalized deep learning denoising strategy for low-count PET images}, \textbf{Physics in Medicine \& Biology}, 2022
\end{itemize}

\textbf{Co-Author}
\begin{itemize}[leftmargin=*]
\item Xia M., Xie H., Liu Q., et al., \textit{LeqMod: adaptable lesion-quantification-consistent modulation for deep learning low-count PET image denoising}, IEEE Transactions on Medical Imaging, 2025
\item Zhou B., Xie H., Liu Q., et al., \textit{FedFTN: personalized federated learning with deep feature transformation network for multi-institutional low-count PET denoising}, Medical Image Analysis, 2023
\end{itemize}

\section{Awards}
\textbf{Society of Nuclear Medicine and Molecular Imaging (SNMMI)}
\begin{itemize}[leftmargin=*]
\item Young Investigator Award (1st Place), 2025
\item Poster Award (2nd Place), 2024
\item Young Investigator Award (2nd Place), 2023
\end{itemize}

\textbf{Mitacs}
\begin{itemize}[leftmargin=*]
\item Globalink Research Award, 2019
\end{itemize}